\section{Introduction}
Before neural dense retrieval, production search and recommendation systems
were built on a two-stage pattern: a cheap lexical or statistical retriever
narrows a large catalogue to a shortlist, and a boosted-tree ranker --
trained on explicit, auditable features -- reorders that shortlist for
quality. This project reproduces that pattern end to end for matching
candidate skill profiles to LinkedIn-style job postings: BM25 lexical
retrieval (stage~1) followed by LambdaMART re-ranking (stage~2) on
handcrafted, point-in-time-safe features. The scope is deliberately
classical: no neural network, no learned embedding space, and no deep
retriever appear anywhere in the stack. The central lesson is that a fully
lexical/statistical pipeline remains a strong, inexpensive, and highly
interpretable baseline, and that taxonomy normalization -- collapsing skill
synonyms to one canonical token -- is the single engineering lever available
to such a system in place of a learned embedding space.
